% =============================================================================
\section{Remediation and Incident Response}
\label{sec:remediation}

\textbf{The fix is the remediation.} Upgrade to Gitea 1.27.1 or later: the
affected range is \texttt{>=1.17, <1.27.1} and 1.27.1 is the first patched
release. Version is verifiable unauthenticated through
\texttt{GET /api/v1/version}, and our precondition measurements
(Section~\ref{sec:precond}) mean that an operator should not spend time
checking whether their configuration is ``bad enough'' --- in the reference
Docker deployment all three vendor preconditions hold by default. Because
CISA added this CVE to the KEV catalog on 2026-08-25 with a federal deadline
of 2026-08-28, and because the advisory itself ships a working exploit, an
internet-reachable Gitea below 1.27.1 should be treated as already targeted,
not hypothetically vulnerable.

\textbf{What an interim mitigation can and cannot do.} Removing open
registration raises the attacker's prerequisite from \emph{none} to
\emph{one ordinary account} --- it does not mitigate the flaw, because the
vulnerability triggers from ordinary write access to any repository
(Section~\ref{sec:precond-cred}). Do not treat registration policy, WAF rules
on patch content, or response-code alerting as substitutes: Sections
\ref{sec:fixchange} and \ref{sec:detect} explain why each of those is weak
against this specific fix shape. If an upgrade must wait, restrict who can
create repositories and who can reach \texttt{diffpatch}, and deploy the
process-lineage detection of Section~\ref{sec:detect} so that an attempt
becomes visible at the only layer that discriminates.

\textbf{Incident response, keyed to measured artifacts.}
\begin{enumerate}
\item \emph{Triage (read-only).} Enumerate running processes for a
\texttt{git} command with a shell descendant, specifically
\texttt{git write-tree} $\rightarrow$ \texttt{/bin/sh\allowbreak
hooks/}\ldots; list files under
\texttt{/data/gitea/tmp/\allowbreak local-repo/\allowbreak upload.git*/hooks/}
(the window is short, so do this continuously rather than on demand); review
reverse-proxy or access logs for repeated \texttt{diffpatch} POSTs to the same
repository inside a one-second window, and for 500s on that route.
\item \emph{Assume credential exposure proportional to the measured impact.}
In our lab the hook read a 1323-byte \texttt{app.ini} whose three secret-bearing
lines are \texttt{SECRET\_KEY} (empty here), \texttt{INTERNAL\_TOKEN} and
\texttt{JWT\_SECRET}, and wrote an arbitrary file under \texttt{/data/gitea}. No
database credential was obtained in this experiment (the lab uses SQLite with an
empty \texttt{PASSWD}), but a readable \texttt{app.ini} is the ordinary path to one
wherever the database has a real password, so if execution is confirmed, rotate
application secrets and database credentials,
and treat repository contents and any integration tokens reachable from that
configuration as exposed --- not because we demonstrated misuse, but because
we demonstrated readability and writability.
\item \emph{Eradication and recovery.} Rebuild rather than clean a host where
arbitrary write by the service user is confirmed inside the configuration
volume. Upgrade first, then verify the fix \emph{as a behaviour}: our A/B
protocol (Section~\ref{sec:fix}) is 22 cycles of an inert canary payload with
an arrival monitor --- the same three lines of evidence (no canary, HTTP 500
on the second submission, payload still observed on disk) are what an operator
should require before declaring an instance fixed. Re-running the vendor PoC
is not a fix test: it proves only that the exploit ``did not appear to work''
on one run.
\end{enumerate}

% =============================================================================
\section{Limitations and Negative Results}
\label{sec:limitations}

This section is the part of the paper an advisory-shaped summary would drop.
It stays in, because four of our own measurements were wrong or unprovable,
and a reader cannot weigh the positive results without them.

\subsection{Methodological failures we made and corrected}
\label{sec:negative}

\begin{enumerate}
\item \textbf{An invalid ``route is absent'' measurement.} Our first check of
the \texttt{diffpatch} precondition was
\texttt{grep -c diffpatch /usr/local/bin/gitea}, which returned 0. That was
meaningless: the file is a \textbf{581-byte} bash wrapper, and the real binary
is \texttt{/app/gitea/gitea} (124\,MB), where the route strings are present.
Had we not checked file size, we would have published a false negative about
our own feasibility gate.
\item \textbf{A timestamp format that silently truncated.} The canary used
\texttt{date -u +\%Y-\%m-\%dT\%H:\%M:\%S.\%6N\%z}; busybox \texttt{date} does
not expand \texttt{\%N}, so \texttt{executed\_at\_utc} ends as
\texttt{2026-08-27T14:33:22.} (visible in Listing~\ref{lst:canary}). Nothing
is invalidated --- the proof is the planted nonce plus the process lineage,
not the timestamp --- but the field is not usable as sub-second evidence.
\item \textbf{A repository count we do not use.} The impact probe reported
\texttt{total\_repos=0}, which contradicts the same probe's own listing of
two repository paths. The cause is an escaping bug in a hook string of ours
(\texttt{``*.git''} inside a Go-quoted command), not a property of the target.
The count is therefore discarded and no repository total appears anywhere in
this paper.
\item \textbf{An A/B probe we threw away.} A structural probe intended to
record bare/non-bare state and index contents atomically per cycle was
unreliable: the temporary repository lives for milliseconds, our bare check
returned \texttt{error}, and several assertions lost a race against cleanup.
It is \textbf{not} used as evidence. The A/B conclusion rests solely on the
three observations of Table~\ref{tab:ab}.
\item \textbf{A capture that did not exist.} The first \texttt{tcpdump} run
produced no file: without \texttt{sudo} the interface capture fails on
\texttt{CAP\_NET\_RAW}. The wire evidence in this paper was recaptured with
elevated capability (\texttt{evidence/wire-attack.pcap}, 33\,400 bytes, 72
packets). We document it because an unrecorded empty capture is exactly the
kind of evidence that would have been fabricated by omission.
\end{enumerate}

\subsection{Scope limits of the positive results}
\begin{itemize}
\item \textbf{Fix semantics on an official image, not a rebuilt binary.} The
patched arm is the vendor's \texttt{gitea/gitea:1.27.1} image. That is the
strongest available comparison for an operator, but it means our A/B cannot
attribute the outcome to the single hunk of Listing~\ref{lst:fix} alone; 1.27.1
contains other changes. We claim the hunk is decisive from code reading, and
behavioural inertness from observation, not from a single-variable binary
diff.
\item \textbf{The vendor's exfiltration channel is not reproduced.} Our PoC
has no output path by design, so we make no claim about the retrieval branch
technique beyond citing the advisory.
\item \textbf{No availability or privilege-escalation claim.} The write proof
was a single file inside \texttt{/data/gitea}, deleted afterwards; we did not
test container escape, service disruption, or escalation beyond
\texttt{uid=1000(git)}.
\item \textbf{Lab network shape.} Attacker position is the bridge gateway of a
zero-egress lab. Latency, TLS termination, reverse-proxy request buffering and
rate limiting of a production deployment are not measured, and the
$0.82$--$0.89$\,s figure is a lab delivery time, not a real-world one.
\item \textbf{Single lab instantiation per version.} One container per arm,
one configuration per arm. The 22/22 figure is repetition of the attack, not
diversity of deployments.
\end{itemize}

% =============================================================================
\section{Ethics, Reproducibility, and Disclosure Position}
\label{sec:ethics}

\textbf{No discovery claim.} CVE-2026-60004, the mechanism, the vulnerable
command line, and a complete working exploit are the vendor's contribution,
published in GHSA-rcr6-4jqh-j84m on 2026-07-28. This paper states that in the
abstract, in Section~\ref{sec:record}, in the companion video material, and in
the artifact README, so that no reader can mistake a reproduction for a
discovery.

\textbf{Consent and blast radius.} Every experiment ran in a Docker lab we
created and destroyed: \texttt{--internal} bridge with verified zero egress,
fictional hostname \texttt{gitea.lab}, no real users, no real data, no
third-party system. Payloads were canary-first by policy: the maximum side
effect of any proof in this paper is one canary file, one probe file (deleted
after reading), and repository objects inside a container we own. No reverse
shell, no destructive action, no denial of service, no social engineering.

\textbf{Reproducibility.} The companion folder is self-contained:
\texttt{lab/run-lab.sh} brings up both arms and \texttt{lab/verify.sh} asserts
the 16 invariants it checks (internal network flag, both versions reported,
three zero-egress probes --- no route to a public address from either arm and no
name resolution on the vulnerable arm --- plus a host-side routing assertion, the
fictional \texttt{*.lab} hostname, git version floor, executable \texttt{/tmp},
real binary size, presence of the \texttt{diffpatch} string) with a pass/fail
exit code;
\texttt{exploit/cve-2026-60004-canary.py} is a stdlib-only reimplementation;
\texttt{exploit/impact-probe.py} produces the impact table;
\texttt{exploit/ab-monitor.sh} produces the arrival evidence that keeps the
negative arm live; \texttt{evidence/} holds the canary and probe transcripts,
the two PoC run logs, and the wire capture; \texttt{designs/} holds the fact
sheet and the fetched fix diff. Every number in this paper is either in
\texttt{designs/FACTS.md} or in one of those files. A reader with Docker can
reproduce Table~\ref{tab:ab} in one afternoon.

\textbf{Operational note on our own pipeline.} The lab, the payloads, the
measurements and this text were produced by an autonomous agent pipeline under
an operator-approved scope memo, including the negative results of
Section~\ref{sec:negative}. We flag this as a limitation as much as a
convenience: an automated pipeline produced a wrong \texttt{grep} measurement,
a broken timestamp, a bogus repository count and an unusable structural probe,
and each was caught only because the protocol required recording failed
attempts. Pipeline output is trustworthy at exactly the strength of its
negative-result discipline.

% =============================================================================
\section{Conclusion}
\label{sec:conclusion}

CVE-2026-60004 is a composition failure, not a coding error: a patch applied
with \texttt{--index}, a three-way fallback that checks a path out, and a bare
sandbox whose root is \texttt{\$GIT\_DIR}, each defensible alone and lethal
together. The vendor did more than disclose it --- they published the exploit.
What a disclosure cannot supply is the rest of the lifecycle, and that is what
we contribute here: execution proven at process level with the parent named
(\texttt{git write-tree}) and the bare-root equivalence observed
(\texttt{GIT\_DIR=.}), preconditions measured rather than assumed and shown to
hold by default in the official image, a fix verified with a live negative
(22/22 cycles with no execution, arrival positively observed in 2 of 22 samples,
HTTP 201 turning into HTTP 500), impact
quantified instead of listed, and a negative-results log that includes four of
our own broken measurements.

The one thing we would ask a reader to keep is this: \emph{read the fix, not
just the advisory}. Changing \texttt{Clone(\dots,true)} to
\texttt{Clone(\dots,false)} neutralises the exploit without making it
detectable --- the same bytes still land on disk. Any defence built on
recognising the attack's content will fire on benign patches (our mode-bit
signature fired on 3 of 3 legitimate cycles that added an executable file) and
stay silent about whether the host was actually exposed. The defence that works is
behavioural: a \texttt{git} process that spawns a shell.

% =============================================================================
\begin{thebibliography}{99}

\bibitem{ghsa}
Gitea Security Team, ``Remote Code Execution via diffpatch Git Hook
Installation --- GHSA-rcr6-4jqh-j84m (CVE-2026-60004), affected
\texttt{>=1.17, <1.27.1}, patched in 1.27.1, published 2026-07-28, including
a complete Python proof-of-concept,''
\url{https://github.com/advisories/GHSA-rcr6-4jqh-j84m}.

\bibitem{nvd}
NIST NVD, ``CVE-2026-60004'' (status Analyzed; published 2026-08-26;
CVSS 3.1 9.8 Critical, secondary assessment; CWE-94; no CVSS v4.0 score),
\url{https://nvd.nist.gov/vuln/detail/CVE-2026-60004}.

\bibitem{fixcommit}
Gitea, ``refactor: git patch apply (\#38637) (\#38638),'' commit
\texttt{d7bc52be\hspace{0pt}eadff4be5f56\hspace{0pt}90de4d5de4\hspace{0pt}2abd10affe},
2026-07-26, \texttt{services/repository/-files/patch.go}.

\bibitem{kev}
CISA, ``Known Exploited Vulnerabilities Catalog --- CVE-2026-60004, added
2026-08-25, FCEB due date 2026-08-28, catalog version \texttt{2026.08.27},
\texttt{dateAdded} and \texttt{dueDate} reproduced verbatim in the shipped
KEV record, \texttt{knownRansomwareCampaignUse =
Unknown}, notes referencing BOD 26-04,''
\url{https://www.cisa.gov/known-exploited-vulnerabilities-catalog}.

\bibitem{bod}
CISA, ``Binding Operational Directive 26-04'' (staggered remediation windows
of 3, 14 and 60 days referenced by the KEV entry for this CVE).

\bibitem{githooks}
J.~Hamano and the Git project, ``git hooks --- \texttt{post-index-change},''
Git documentation, \url{https://git-scm.com/docs/githooks}.

\bibitem{gitapply}
J.~Hamano and the Git project, ``git-apply --- \texttt{--index},
\texttt{-\/-cached}, \texttt{-3},'' Git documentation,
\url{https://git-scm.com/docs/git-apply}.

\bibitem{cwe94}
MITRE, ``CWE-94: Improper Control of Generation of Code ('Code
Injection'),'' \url{https://cwe.mitre.org/data/definitions/94.html}.

\bibitem{attack}
MITRE, ``MITRE ATT\&CK: T1190, T1059.004, T1083, T1552.001, T1546.004,''
\url{https://attack.mitre.org/}.

\end{thebibliography}

\end{document}
